\section*{Hoofdstuk \ref{ch:b1:comparative-planetology} --- Vergelijkende planetologie}

\begin{solution}{exo:b1:comparative-planetology:1}
$V = \tfrac43\pi(\qty{6.9911e7}{m})^3 = \qty{1.4313e24}{m^3}$ en
$\bar\rho = \num{1898.13e24}/\num{1.4313e24} = \qty{1326}{kg.m^{-3}}$, slechts
1.33 keer de dichtheid van water, hoewel Jupiter 318 keer zo zwaar is als de aarde: zij
bestaat vooral uit de lichtste elementen.
\end{solution}

\begin{solution}{exo:b1:comparative-planetology:2}
Mercurius (straal \qty{2440}{km}, \qty{5429}{kg.m^{-3}}) en Venus
(\qty{6052}{km}, \qty{5243}{kg.m^{-3}}) zijn klein en dichter dan gesteente:
\omterm{def:b1:comparative-planetology:planet-types}{aardachtige planeten}, van gesteente en ijzer. Saturnus (\qty{60268}{km},
\qty{687}{kg.m^{-3}}) is enorm en lichter dan water: een \omterm{def:b1:comparative-planetology:planet-types}{gasreus}, vooral uit
waterstof en helium. Uranus (\qty{25559}{km}, \qty{1270}{kg.m^{-3}}) is een reus, maar
veel kleiner dan Saturnus en bijna twee keer zo dicht: een \omterm{def:b1:comparative-planetology:planet-types}{ijsreus}, waarvan het water, het
methaan en de ammoniak zwaarder zijn dan waterstof.
\end{solution}

\begin{solution}{exo:b1:comparative-planetology:3}
$T_\mathrm{eq} = (2601.3\times0.23/(4\times\num{5.6704e-8}))^{1/4} =
(\num{2.638e9})^{1/4} = \qty{226.6}{K}$. Venus ontvangt 1.9 keer meer zonlicht dan de aarde,
maar haar wolken kaatsen \qty{77}{\percent} ervan terug: zij absorbeert slechts
$2601.3\times0.23 = \qty{598}{W.m^{-2}}$, tegen $1361\times0.706 = \qty{961}{W.m^{-2}}$ voor de aarde.
\end{solution}

\begin{solution}{exo:b1:comparative-planetology:4}
Van het meest homogene naar het meest geconcentreerde: de maan, Mars, de aarde, Jupiter,
Saturnus. Een puntmassa in het middelpunt heeft $I = 0$, dus $C = 0$.
\end{solution}

\begin{solution}{exo:b1:comparative-planetology:5}
$T = 15.945421\times86400 = \qty{1377684}{s}$;
$GM = 4\pi^2(\num{1.22187e9})^3/\num{1377684}^2 = \qty{3.794e16}{m^3.s^{-2}}$;
$M = \num{3.794e16}/\num{6.6743e-11} = \qty{5.685e26}{kg}$, binnen
\qty{0.03}{\percent} van \qty{568.32e24}{kg}.
\end{solution}

\begin{solution}{exo:b1:comparative-planetology:6}
$S = 1361\times(149.6/228.0)^2 = \qty{585.9}{W.m^{-2}}$ bij Mars en
$1361\times(149.6/778.5)^2 = \qty{50.26}{W.m^{-2}}$ bij Jupiter, in overeenstemming
met de tabellen (het kleine verschil bij Mars komt door de afronding van de gemiddelde afstanden).
\end{solution}

\begin{solution}{exo:b1:comparative-planetology:7}
$S \propto d^{-2}$ en $T_\mathrm{eq} \propto S^{1/4}$, dus $T_\mathrm{eq}
\propto d^{-1/2}$. Voor Jupiter: $T_\mathrm{eq} = (50.26\times0.657/(4\times
\num{5.6704e-8}))^{1/4} = \qty{109.8}{K}$. Een planeet die meer uitstraalt dan zij
absorbeert, moet een inwendige warmtebron hebben: Jupiter zou warmte afgeven die sinds haar vorming
is opgeslagen, of die wordt opgewekt terwijl zij langzaam krimpt.
\end{solution}

\begin{solution}{exo:b1:comparative-planetology:8}
Venus: $737 - 226.6 = \qty{510}{K}$, \qty{225}{\percent} van $T_\mathrm{eq}$.
De aarde: $T_\mathrm{eq} = (1361\times0.706/(4\sigma))^{1/4} = \qty{255.1}{K}$,
$288 - 255.1 = \qty{33}{K}$, \qty{13}{\percent}. Mars: $214 - 209.8 =
\qty{4}{K}$, \qty{2}{\percent}.
\end{solution}

\begin{solution}{exo:b1:comparative-planetology:9}
Met $\mathrm dm = 4\pi\rho r^2\,\mathrm dr$: $I = \int_0^R\tfrac23r^2\,4\pi
\rho r^2\,\mathrm dr = \tfrac{8\pi}{15}\rho R^5$, en $MR^2 = \tfrac43\pi\rho
R^5$, dus $C = \tfrac{8}{15}\cdot\tfrac34 = \tfrac25$. Een holle bolschil heeft al haar massa
op $r = R$: $I = \tfrac23MR^2$, $C = \tfrac23$.
\end{solution}

\begin{solution}{exo:b1:comparative-planetology:10}
$\tfrac52C\bar\rho = 2.5\times0.35\times5429 = 4750.4$;
$x^2 = (4750.4 - 3300)/(5429 - 3300) = 1450.4/2129 = 0.6812$, $x = 0.8254$,
kernstraal $0.8254\times2439.7 = \qty{2014}{km}$. $x^3 = 0.5623$ en
$\rho_c = 3300 + 2129/0.5623 = \qty{7086}{kg.m^{-3}}$. De kern bevat
$0.5623\times7086/5429 = \qty{73}{\percent}$ van de massa. De kern van Mercurius
bereikt \qty{83}{\percent} van haar straal, tegen $3485/6371 =
\qty{55}{\percent}$ voor de aarde: Mercurius bestaat grotendeels uit kern, met een dunne
rotsachtige mantel.
\end{solution}

\begin{solution}{exo:b1:comparative-planetology:11}
(a) De functie $t\mapsto t^4$ is convex, dus
$\bigl(\tfrac{T_1 + T_2}{2}\bigr)^4 \le \tfrac{T_1^4 + T_2^4}{2} =
T_\mathrm{eq}^4$, met gelijkheid alleen als $T_1 = T_2$. (b) $T_1^4 =
2\times270^4 - 100^4 = \num{1.0529e10}$, $T_1 = \qty{320}{K}$, en de gemiddelde
temperatuur is $(320 + 100)/2 = \qty{210}{K}$, \qty{60}{K} onder
$T_\mathrm{eq}$.
\end{solution}

\begin{solution}{exo:b1:comparative-planetology:12}
$m = 4\pi R^2p/g = 4\pi\times(\num{6.0518e6})^2\times\num{9.2e6}/8.9 =
\num{4.602e14}\times\num{1.034e6} = \qty{4.76e20}{kg}$: 93 keer de massa van de
atmosfeer van de aarde, en \num{9.8e-5} van de massa van Venus. De aarde houdt een vergelijkbare
hoeveelheid koolstof vast, niet in haar lucht maar in carbonaatgesteenten (kalksteen) en in de
oceaan, waar vloeibaar water de opslag mogelijk maakte.
\end{solution}

\begin{solution}{pb:b1:comparative-planetology:1}
\textbf{1.} $T = 0.31891\times86400 = \qty{27553.8}{s}$.

\textbf{2.} $GM = 4\pi^2\times(\num{9.378e6})^3/\num{27553.8}^2 =
\qty{4.289e13}{m^3.s^{-2}}$.

\textbf{3.} $M = \num{4.289e13}/\num{6.6743e-11} = \qty{6.426e23}{kg}$,
binnen \qty{0.14}{\percent} van de getabelleerde \qty{6.4169e23}{kg}.

\textbf{4.} $\num{1.06e16}/\num{6.426e23} = \num{1.6e-8}$: Fobos is
verwaarloosbaar.

\textbf{5.} $v = 2\pi a/T = 2\pi\times\num{9.378e6}/27553.8 =
\qty{2138}{m.s^{-1}}$, ongeveer \qty{2.1}{km.s^{-1}}.

\textbf{6.} $T = \qty{7.65}{h}$, 3.2 keer korter dan de rotatie van Mars: Fobos draait sneller
rond dan de grond onder haar draait, dus zij haalt de rotatie in. Gezien vanaf Mars komt zij in het
westen op en in het oosten onder.

\textbf{7.} $V = \tfrac43\pi(\num{3.3895e6})^3 = \qty{1.6312e20}{m^3}$.

\textbf{8.} $\bar\rho = \num{6.426e23}/\num{1.6312e20} = 3939$, dat is
\qty{3940}{kg.m^{-3}}.

\textbf{9.} $3940/5514 = 0.715$. De aarde, groter en zwaarder, perst haar eigen binnenste sterker
samen; en haar ijzeren kern is een groter deel van de planeet.

\textbf{10.} $g = GM/R^2 = \num{4.289e13}/(\num{3.3895e6})^2 =
\qty{3.73}{m.s^{-2}}$, \qty{38}{\percent} van die van de aarde.

\textbf{11.} $v_\mathrm{esc} = \sqrt{2GM/R} = \sqrt{2\times\num{4.289e13}/
\num{3.3895e6}} = \qty{5.03}{km.s^{-1}}$.

\textbf{12.} $T_\mathrm{eq} = (586.2\times0.75/(4\times\num{5.6704e-8}))^{1/4}
= \qty{209.8}{K}$.

\textbf{13.} $214 - 209.8 = \qty{4}{K}$: een zwak broeikasopwarmingseffect.

\textbf{14.} $m = p/g = 636/3.733 = \qty{170}{kg.m^{-2}}$; voor de aarde
$101400/9.82 = \qty{10326}{kg.m^{-2}}$, 61 keer meer.

\textbf{15.} $4\pi R^2\times170.4 = \num{1.444e14}\times170.4 =
\qty{2.46e16}{kg}$.

\textbf{16.} Boven elke vierkante meter ligt zestig keer minder gas dan op de aarde, dus de
infraroodabsorptie, hoewel door koolstofdioxide, is zwak; en de lucht van Mars bevat zeer weinig
waterdamp, de andere belangrijke absorbeerder.

\textbf{17.} Een homogeen Mars zou $C = 0.4$ hebben; de gemeten 0.366 toont aan dat haar dichtheid
naar het midden toe toeneemt: Mars heeft een dichte kern.

\textbf{18.} $\bar\rho = \rho_m + (\rho_c - \rho_m)x^3$.

\textbf{19.} $C = \tfrac25[\rho_m + (\rho_c - \rho_m)x^5]/\bar\rho$.

\textbf{20.} Uit 18 en 19 volgt $(\rho_c - \rho_m)x^3 = \bar\rho - \rho_m$ en
$(\rho_c - \rho_m)x^5 = \tfrac52C\bar\rho - \rho_m$; delen geeft $x^2$.

\textbf{21.} $\tfrac52C\bar\rho = 2.5\times0.366\times3940 = 3605.1$;
$x^2 = (3605.1 - 3500)/(3940 - 3500) = 105.1/440 = 0.2389$, $x = 0.4887$.
$x^3 = 0.1167$ en $\rho_c = 3500 + 440/0.1167 = \qty{7269}{kg.m^{-3}}$.

\textbf{22.} $x^3\rho_c/\bar\rho = 0.1167\times7269/3940 = 0.215$, dat is
\qty{21.5}{\percent} van de massa.

\textbf{23.} Met $\rho_m = 3450$: $x^2 = 155.1/490 = 0.3165$, $x = 0.5626$,
kernstraal \qty{1907}{km}, $\rho_c = \qty{6202}{kg.m^{-3}}$. Met $\rho_m =
3550$: $x^2 = 55.1/390 = 0.1413$, $x = 0.3759$, \qty{1274}{km},
$\rho_c = \qty{10894}{kg.m^{-3}}$. Een verandering van \qty{1.4}{\percent} in de
aangenomen dichtheid van de mantel verschuift de kernstraal met enkele honderden kilometers:
de gemiddelde dichtheid en $C$ alleen bepalen de kern van Mars niet nauwkeurig, en daarom zijn
seismometers erheen gestuurd.

\textbf{24.} $0.4887\times3389.5 = \qty{1657}{km}$: de straal van de kern
van Mars volgens het model is \textbf{ongeveer \qty[text-series-to-math=true]{1660}{km}}, binnen het
recente seismische bereik van \qty{1650}{km} tot \qty{1675}{km} en duidelijk onder de
\qty{1830}{km} van de planetaire tabel --- een overeenstemming die vraag 23 laat zien deels
geluk te zijn bij de keuze van $\rho_m$.
\end{solution}
